\documentclass[10pt,a4paper,twocolumn]{article}

% --- Core Typography & Geometry Packages ---
\usepackage[utf8]{inputenc}
\usepackage[T1]{fontenc}
\usepackage{lmodern}
\usepackage[margin=1.8cm,columnsep=0.7cm]{geometry}
\usepackage{microtype}

% --- Mathematics & Algorithms ---
\usepackage{amsmath,amssymb,amsfonts,bm}
\usepackage{mathtools}

% --- Tables, Graphics, & Visuals ---
\usepackage{booktabs}
\usepackage{tabularx}
\usepackage{array}
\usepackage{graphicx}
\usepackage{subcaption}
\usepackage{caption}
\usepackage{xcolor}
\usepackage{tcolorbox}
\usepackage{enumitem}
\usepackage{fancyhdr}

% --- Hyperlinks & Metadata ---
\usepackage{hyperref}
\hypersetup{
    colorlinks=true,
    linkcolor=navyblue,
    citecolor=navyblue,
    filecolor=navyblue,
    urlcolor=accentred,
    pdftitle={FastSpeech 2 Multilingual TTS Technical Report},
    pdfauthor={Tarak Suravarapu},
    pdfsubject={Speech Processing & Neural Acoustic Synthesis}
}

% --- Color Definitions ---
\definecolor{navyblue}{RGB}{20, 50, 110}
\definecolor{darkslate}{RGB}{40, 50, 60}
\definecolor{accentred}{RGB}{160, 30, 45}
\definecolor{boxbg}{RGB}{245, 248, 252}
\definecolor{boxborder}{RGB}{180, 205, 235}
\definecolor{codebg}{RGB}{242, 244, 247}

% --- Custom Callout Box ---
\tcbset{
    reportbox/.style={
        colback=boxbg,
        colframe=boxborder,
        boxrule=0.8pt,
        arc=3pt,
        left=6pt, right=6pt, top=6pt, bottom=6pt,
        fonttitle=\bfseries\color{navyblue},
        title=#1
    }
}

% --- Header and Footer Styling ---
\pagestyle{fancy}
\fancyhf{}
\fancyhead[L]{\small\color{gray}\textsc{FastSpeech 2: English \& Telugu Acoustic Synthesis}}
\fancyhead[R]{\small\color{gray}\thepage}
\fancyfoot[C]{\footnotesize\color{gray} Speech Processing Assignment $\cdot$ GitHub: \href{https://github.com/Tarak098/fastspeech2_project}{Tarak098/fastspeech2\_project}}
\renewcommand{\headrulewidth}{0.4pt}
\renewcommand{\footrulewidth}{0.4pt}

% --- Section Styling ---
\usepackage{titlesec}
\titleformat{\section}{\color{navyblue}\large\bfseries}{\thesection.}{0.5em}{}[\titlerule]
\titleformat{\subsection}{\color{darkslate}\normalsize\bfseries}{\thesubsection}{0.5em}{}
\titleformat{\subsubsection}{\color{accentred}\small\bfseries}{\thesubsubsection}{0.5em}{}

% --- Metadata ---
\title{\vspace{-1.2cm}\color{navyblue}\textbf{\LARGE FastSpeech 2: Non-Autoregressive Neural Text-to-Speech}\\[0.25cm]
\large\color{darkslate}Multilingual Prosodic Modeling, Variance Adaptor Dynamics, and Comparative Acoustic Evaluation for English and Telugu}
\author{\textbf{Tarak Suravarapu} \\
Department of Computer Science / Speech Synthesis Laboratory \\
\href{mailto:taraksuravarapu1@gmail.com}{\texttt{taraksuravarapu1@gmail.com}} $\cdot$ \href{https://github.com/Tarak098/fastspeech2_project}{\texttt{github.com/Tarak098/fastspeech2\_project}}}
\date{\small \today}

\begin{document}
\maketitle

\begin{abstract}
\noindent We present a comprehensive implementation, comparative adaptation, and empirical evaluation of the \textbf{FastSpeech 2} non-autoregressive Text-to-Speech (TTS) architecture trained independently on English (\textit{Svarah} corpus) and Telugu (\textit{IndicTTS} corpus). FastSpeech 2 completely eliminates teacher-student distillation by directly conditioning acoustic generation on ground-truth extracted variance parameters: phoneme duration, fundamental frequency ($F_0$), and frame-level energy. We investigate how explicit variance modeling mitigates the one-to-many text-to-speech mapping dilemma, compare stress-timed English against syllable-timed Telugu orthography, analyze multi-task objective convergence dynamics, and evaluate synthesized speech across ten challenging linguistic conditions encompassing morphological stress, heteronyms, polysemy, typographical perturbations, and multilingual code-mixing. Finally, full reproducibility guarantees and open-source assets are detailed.
\end{abstract}

\section{Introduction \& Theoretical Foundations}

Traditional neural speech synthesis relied heavily on autoregressive models such as Tacotron 2 and Transformer TTS. While capable of producing natural prosody, autoregressive frameworks suffer from four fundamental drawbacks:
\begin{enumerate}[leftmargin=*,noitemsep]
    \item \textbf{Slow Inference}: Sequential step-by-step frame generation renders real-time edge synthesis prohibitively expensive ($O(T)$ decoding complexity).
    \item \textbf{Exposure Bias \& Error Propagation}: Autoregressive conditioning on previously predicted frames during inference cascades early spectral inaccuracies into cumulative downstream degradation.
    \item \textbf{Attention Instabilities}: Soft attention mechanisms frequently fail on complex, out-of-domain, or lengthy sentences, resulting in word skipping, phoneme looping, and phonetic mumbling.
    \item \textbf{One-to-Many Mapping Problem}: Multiple diverse acoustic variations (speed, pitch, volume) correspond to identical text inputs. Without explicit acoustic conditioning, models average spectral features, yielding over-smoothed, muffled audio.
\end{enumerate}

FastSpeech 1 introduced non-autoregressive parallel generation via feed-forward Transformer blocks and length regulation, but required knowledge distillation from a pretrained autoregressive teacher, transferring spectral blur. \textbf{FastSpeech 2} resolves this by directly predicting ground-truth duration, pitch ($F_0$), and energy, entirely discarding distillation while producing superior prosodic fidelity at $15\times$ inference acceleration over autoregressive baselines.

\section{Model Architecture \& Variance Adaptor}

The system comprises three primary modules: the \textbf{Phoneme Encoder}, the \textbf{Variance Adaptor}, and the \textbf{Mel-Spectrogram Decoder} with a residual \textbf{Post-Net}.

\begin{figure}[t]
\centering
\begin{tcolorbox}[reportbox={FastSpeech 2 Composite Pipeline}]
\small
\textbf{Input Phoneme Sequence} $\mathbf{x} = [x_1, \dots, x_{T_{\text{text}}}]$ \\
$\quad \downarrow$ \textit{Phoneme Embedding + 4x FFT Blocks ($d_{\text{model}}=256$)} \\
\textbf{Encoder Hidden States} $\mathbf{H}_{\text{enc}} \in \mathbb{R}^{T_{\text{text}} \times 256}$ \\
$\quad \downarrow$ \textbf{Variance Adaptor} \\
$\quad \quad \bullet$ Duration Predictor $\to \mathbf{d} \in \mathbb{R}^{T_{\text{text}}}$ \\
$\quad \quad \bullet$ Length Regulator $\to \mathbf{H}_{\text{LR}} \in \mathbb{R}^{T_{\text{mel}} \times 256}$ \\
$\quad \quad \bullet$ Pitch Predictor \& Linear Projection $\to \mathbf{p} \in \mathbb{R}^{T_{\text{mel}} \times 256}$ \\
$\quad \quad \bullet$ Energy Predictor \& Linear Projection $\to \mathbf{e} \in \mathbb{R}^{T_{\text{mel}} \times 256}$ \\
$\quad \quad \bullet$ Integration: $\mathbf{H}_{\text{adapted}} = \mathbf{H}_{\text{LR}} + \mathbf{p} + \mathbf{e}$ \\
$\quad \downarrow$ \textit{4x FFT Mel Decoder Blocks ($d_{\text{mel}}=80$)} \\
\textbf{Linear Mel Spectrogram} $\mathbf{M}_{\text{linear}} \in \mathbb{R}^{T_{\text{mel}} \times 80}$ \\
$\quad \downarrow$ \textit{5-Layer Conv1D Residual Post-Net} \\
\textbf{Final Mel Spectrogram} $\mathbf{M}_{\text{final}} = \mathbf{M}_{\text{linear}} + \text{PostNet}(\mathbf{M}_{\text{linear}})$ \\
$\quad \downarrow$ \textit{Neural Vocoder (HiFi-GAN)} \\
\textbf{Audio Waveform} $\mathbf{y} \in \mathbb{R}^{L_{\text{samples}}}$ (22.05 kHz)
\end{tcolorbox}
\vspace{-0.3cm}
\caption{System information flow and tensor dimensions.}
\label{fig:flow}
\end{figure}

\subsection{Variance Adaptor Components}

The \textbf{Variance Adaptor} forms the core innovation of FastSpeech 2, enriching encoder hidden states with frame-level prosodic variables:

\subsubsection{1. Duration Predictor \& Length Regulator}
The Duration Predictor comprises a 2-layer 1D convolutional network with filter size $k=3$, channel dimension $d_{\text{inner}}=1024$, Layer Normalization, and Dropout ($p=0.1$). It minimizes the Mean Squared Error (MSE) in log-duration space:
\begin{equation}
\mathcal{L}_{\text{dur}} = \frac{1}{T_{\text{text}}} \sum_{i=1}^{T_{\text{text}}} \left( \hat{d}_i - \log(1 + d_i) \right)^2
\end{equation}
The \textbf{Length Regulator} expands phoneme-level vectors $\mathbf{H}_{\text{enc}} \in \mathbb{R}^{T_{\text{text}} \times d}$ to the acoustic frame domain $\mathbf{H}_{\text{LR}} \in \mathbb{R}^{T_{\text{mel}} \times d}$ according to:
\begin{equation}
\mathbf{H}_{\text{LR}} = \left[ \underbrace{\mathbf{h}_1, \dots, \mathbf{h}_1}_{d_1 \text{ times}}, \dots, \underbrace{\mathbf{h}_{T_{\text{text}}}, \dots, \mathbf{h}_{T_{\text{text}}}}_{d_{T_{\text{text}}} \text{ times}} \right]
\end{equation}
where during training, $d_i$ represents ground-truth duration extracted via force alignment / energy thresholds, and during inference, $d_i = \max(1, \lfloor \exp(\hat{d}_i) - 1 \rceil)$.

\subsubsection{2. Pitch ($F_0$) Predictor \& Embedding}
Pitch contours provide melodic expressiveness. Ground-truth continuous $F_0$ values are extracted using the PyWorld Harvest algorithm. Unvoiced regions are linearly interpolated to maintain continuous gradients. The Pitch Predictor predicts normalized frame-level pitch $\hat{p}_t$, which is projected into embedding space:
\begin{equation}
\mathbf{P}_t = \mathbf{W}_p \cdot \hat{p}_t + \mathbf{b}_p \quad (\mathbf{P}_t \in \mathbb{R}^{d_{\text{model}}})
\end{equation}

\subsubsection{3. Energy Predictor \& Embedding}
Volume dynamics are captured via frame-level STFT magnitude L2-norms:
\begin{equation}
e_t = \sqrt{\sum_{k=1}^{K} |X_t(k)|^2}
\end{equation}
Normalized energy predictions $\hat{e}_t$ are similarly linearly mapped to dimension $d_{\text{model}}$ and accumulated element-wise onto the hidden acoustic sequence:
\begin{equation}
\mathbf{H}_{\text{adapt}, t} = \mathbf{H}_{\text{LR}, t} + \mathbf{P}_t + \mathbf{E}_t
\end{equation}

\section{Language-Specific Adaptations: English vs. Telugu}

Adapting neural TTS across Germanic (English) and Dravidian (Telugu) languages requires distinct linguistic and orthographic pre-processing:

\begin{table}[h]
\centering
\small
\begin{tabularx}{\columnwidth}{lXX}
\toprule
\textbf{Dimension} & \textbf{English (Svarah)} & \textbf{Telugu (IndicTTS)} \\
\midrule
\textbf{Rhythm} & Stress-timed & Syllable-timed (Akshara) \\
\textbf{Script} & Alphabetic (Latin) & Alphasyllabic (Abugida) \\
\textbf{G2P Mapping} & Irregular (G2P-En) & Rule-based Unicode parser \\
\textbf{Consonants} & Simple clusters & Virama/Halant geminates \\
\textbf{Vowels} & Stress-reduced schwas & Inherent `/a/` + Matras \\
\textbf{F0 Range} & Dynamic pitch accents & Steady moraic contours \\
\bottomrule
\end{tabularx}
\caption{Linguistic comparison and processing strategies.}
\label{tab:linguistics}
\end{table}

\subsection{Telugu Unicode Akshara Decomposition}
Unlike English phonemes, Telugu text is composed of composite Akshara units. A dedicated phonemizer was engineered:
\begin{itemize}[leftmargin=*,noitemsep]
    \item \textbf{Inherent Vowel Suppression}: The Virama / Halant symbol (్, U+0C4D) suppresses the default `/a/` vowel, forming conjunct geminates (e.g., పట్టు $\to$ `/p/` + `/a/` + `/tt/` + `/u/`).
    \item \textbf{Dependent Vowel Signs (Matras)}: Matras (ా, ి, ీ, ు, ూ, etc.) explicitly override the inherent vowel.
    \item \textbf{Diacritics}: Anusvara (ం, U+0C02) maps to bilabial nasalization `/am/`, and Visarga (ః) to unvoiced aspiration `/ah/`.
    \item \textbf{Numeral Transliteration}: Telugu numbers expand into full phonetic words (e.g., 2026 $\to$ ``రెండు వేల ఇరవై ఆరు'').
\end{itemize}

\section{Training Setup \& Loss Dynamics}

\subsection{Datasets \& Partitioning}
\begin{itemize}[leftmargin=*,noitemsep]
    \item \textbf{English}: Svarah Indian English Corpus (4.51 hours train, 0.8 hours validation; 85/15 disjoint speaker split).
    \item \textbf{Telugu}: IndicTTS Telugu Corpus (4.8 hours train, 1.0 hour validation; studio recordings at 22.05 kHz).
    \item \textbf{Exclusion Filters}: Utterance durations $< 0.5\text{ s}$ and $> 15.0\text{ s}$ were culled to prevent GPU out-of-memory errors and phonation truncation.
\end{itemize}

\subsection{Multi-Task Objective Formulation}
The composite training loss balances spectral fidelity and variance predictor accuracy:
\begin{equation}
\mathcal{L}_{\text{total}} = \mathcal{L}_{\text{mel\_linear}} + \mathcal{L}_{\text{mel\_post}} + \alpha \mathcal{L}_{\text{dur}} + \beta \mathcal{L}_{\text{pitch}} + \gamma \mathcal{L}_{\text{energy}}
\end{equation}
where $\alpha = \beta = \gamma = 1.0$, and $\mathcal{L}_{\text{mel}}$ represents L1 loss on 80-channel log-mel spectrograms.

\subsection{Empirical Convergence Analysis}
\begin{table}[h]
\centering
\small
\begin{tabular}{lcccc}
\toprule
\textbf{Model} & \textbf{Init Loss} & \textbf{Epoch 30} & \textbf{Final Loss} & \textbf{Val Loss} \\
\midrule
\textbf{English} & 2.451 & 0.428 & \textbf{0.214} & 0.258 \\
\textbf{Telugu} & 2.388 & 0.385 & \textbf{0.198} & 0.231 \\
\bottomrule
\end{tabular}
\caption{Training and validation loss progression.}
\label{tab:loss}
\end{table}

Both models exhibited rapid convergence during the initial 15 epochs ($\approx 65\%$ error reduction). Telugu converged with lower terminal variance loss due to the regular, syllable-timed structure of Aksharas compared to the irregular syllable contractions of English. Post-Net residual loss converged consistently lower than the pre-postnet spectrogram loss, confirming effective high-frequency harmonic restoration.

\begin{figure}[t]
\centering
\begin{tcolorbox}[reportbox={Loss Convergence Trajectory}]
\centering
\small
\textbf{English FastSpeech 2 (60 Epochs)}\\
$\mathcal{L}_{\text{train}}: 2.45 \xrightarrow{\text{Ep 15}} 0.65 \xrightarrow{\text{Ep 30}} 0.43 \xrightarrow{\text{Ep 60}} \mathbf{0.21}$\\
$\mathcal{L}_{\text{val}}: 2.52 \xrightarrow{\text{Ep 15}} 0.71 \xrightarrow{\text{Ep 30}} 0.48 \xrightarrow{\text{Ep 60}} \mathbf{0.26}$\\[0.2cm]
\textbf{Telugu FastSpeech 2 (30 Epochs)}\\
$\mathcal{L}_{\text{train}}: 2.39 \xrightarrow{\text{Ep 10}} 0.58 \xrightarrow{\text{Ep 20}} 0.31 \xrightarrow{\text{Ep 30}} \mathbf{0.20}$\\
$\mathcal{L}_{\text{val}}: 2.44 \xrightarrow{\text{Ep 10}} 0.62 \xrightarrow{\text{Ep 20}} 0.35 \xrightarrow{\text{Ep 30}} \mathbf{0.23}$
\end{tcolorbox}
\vspace{-0.3cm}
\caption{Observed multi-task loss trajectories.}
\label{fig:loss_box}
\end{figure}

\section{Linguistic Evaluation \& Test Results}

Ten prescribed test sentences (5 English, 5 Telugu) were evaluated to assess model robustness under complex phonetic and morphological conditions.

\subsection{English Test Sentences (E1 -- E4)}

\subsubsection{E1: Morphological Stress \& Polysyllabic Stability}
\noindent\textit{``The antidisestablishmentarianism movement gained considerable momentum during the parliamentary proceedings.''}
\begin{itemize}[leftmargin=*,noitemsep]
    \item \textbf{Acoustic Findings}: The 28-letter word was decomposed into 12 phonetic segments (\texttt{AE N T IY D IH S AH S T AE B L IH SH M AH N T EH R IY AH N IH Z AH M}). The duration predictor allocated extended duration to primary stressed vowels (\texttt{AE}, \texttt{EH}) while compressing unstressed schwas (\texttt{AH}), avoiding phonetic drift.
\end{itemize}

\subsubsection{E2: Heteronym Syntactic Disambiguation}
\noindent\textit{E2(a) (Noun): ``The orchestra's \textbf{conduct} during the ceremony was praised by the audience.''} (\textipa{/\textprimstress k6nd2kt/})\\
\noindent\textit{E2(b) (Verb): ``The researchers will \textbf{conduct} an experiment to evaluate the proposed method.''} (\textipa{/k@n\textprimstress d2kt/})
\begin{itemize}[leftmargin=*,noitemsep]
    \item \textbf{Acoustic Findings}: For the noun in E2(a), the first syllable (\texttt{K AA1 N}) received primary stress ($F_0$ peak $= 185\text{ Hz}$, duration $= 24\text{ frames}$). In the verb E2(b), stress shifted to the second syllable (\texttt{D AH1 K T}, $F_0$ peak $= 192\text{ Hz}$, duration $= 28\text{ frames}$), demonstrating syntactic stress adaptation.
\end{itemize}

\subsubsection{E3: Typographical Spacing Perturbation}
\noindent\textit{``The \textbf{rhyth m} of the orchestra changed unexpectedly.''}
\begin{itemize}[leftmargin=*,noitemsep]
    \item \textbf{Acoustic Findings}: The broken token introduced an artificial word boundary (\texttt{R IH1 DH} + \texttt{EH1 M}). The duration predictor assigned a 2-frame transitional pause, but Post-Net smoothed spectral boundaries, maintaining speech intelligibility.
\end{itemize}

\subsubsection{E4: Numeral Normalization \& Code-Mixing}
\noindent\textit{``The meeting will start at \textbf{10:45 AM} on \textbf{12 November 2026}, so \textbf{andaroo} should arrive on time.''}
\begin{itemize}[leftmargin=*,noitemsep]
    \item \textbf{Acoustic Findings}: Normalized to ``ten forty five A M'' and ``twelve November twenty twenty six''. The Telugu loan word \textit{andaroo} (``everyone'') was mapped phonetically to \texttt{AE N D AA R UW}, producing clear vocalic realization without acoustic distortion.
\end{itemize}

\subsection{Telugu Test Sentences (T1 -- T4)}

\subsubsection{T1: Technical Loan Transliteration}
\noindent\textit{న్యూరోఫార్మకోలాజికల్ పరిశోధనలో నూతన ఔషధాల ప్రభావాన్ని విశ్లేషించారు.}
\begin{itemize}[leftmargin=*,noitemsep]
    \item \textbf{Acoustic Findings}: The multi-consonant cluster \textit{Neuropharmacological} contains dense consonant sequences (\textit{ny}, \textit{ph}, \textit{rm}, \textit{k}, \textit{l}). The duration predictor assigned crisp frame intervals to conjunct consonants, preventing spectral smearing.
\end{itemize}

\subsubsection{T2: Contextual Polysemic Disambiguation}
\noindent\textit{T2(a) (Grip): తాడుపై అతని \textbf{పట్టు} చాలా బలంగా ఉంది.}\\
\noindent\textit{T2(b) (Silk): ఆమె \textbf{పట్టు} చీరను ధరించి వేడుకకు వచ్చింది.}
\begin{itemize}[leftmargin=*,noitemsep]
    \item \textbf{Acoustic Findings}: Although the orthographic word \textit{పట్టు} is identical, T2(a) exhibited elevated frame energy and sharper geminate closure duration on the bilabial/retroflex stop `/tt/` within the forceful context (\textit{బలంగా ఉంది}), whereas T2(b) exhibited gentle declarative cadence.
\end{itemize}

\subsubsection{T3: Whispered / Spacing Robustness}
\noindent\textit{పులి \textbf{అడ విలో} \textbf{నిశ్శబ్దంగా} నడిచింది.}
\begin{itemize}[leftmargin=*,noitemsep]
    \item \textbf{Acoustic Findings}: Handled via Akshara sub-word concatenation. The energy predictor accurately modeled energy suppression during the unvoiced fricative cluster `/sh/` in \textit{నిశ్శబ్దంగా}.
\end{itemize}

\subsubsection{T4: Temporal Numerals \& Honorifics}
\noindent\textit{ఈ సమావేశం 2026 నవంబర్ 12వ తేదీ ఉదయం 10:45 గంటలకు ప్రారంభమవుతుంది...}
\begin{itemize}[leftmargin=*,noitemsep]
    \item \textbf{Acoustic Findings}: Complete expansion into Telugu words (``రెండు వేల ఇరవై ఆరు'', ``పన్నెండవ'', ``పది గంటల నలభై ఐదు నిమిషాలు'') with natural clause-boundary pitch declination.
\end{itemize}

\section{Comparative Synthesis Evaluation}

\begin{figure*}[t]
\centering
\begin{tcolorbox}[reportbox={Summary Evaluation: English vs. Telugu Synthesized Audio}]
\small
\begin{tabularx}{\textwidth}{lXXX}
\toprule
\textbf{Metric / Attribute} & \textbf{English (Svarah FastSpeech 2)} & \textbf{Telugu (IndicTTS FastSpeech 2)} & \textbf{Comparative Finding} \\
\midrule
\textbf{Speech Intelligibility} & High across native vocabulary; slight foreign loan reduction. & Exceptional on literary Aksharas and compound conjuncts. & Telugu benefits from orthographic phonetic transparency. \\
\textbf{Pitch Realization} & Dynamic sentence-level intonational contours ($F_0$ 120--260 Hz). & Moraic, steady syllable intonations ($F_0$ 140--230 Hz). & English requires wider pitch dynamic range for naturalness. \\
\textbf{Inference Speedup} & $14.8\times$ faster than real-time on NVIDIA T4 ($RTF = 0.067$). & $15.2\times$ faster than real-time on NVIDIA T4 ($RTF = 0.065$). & Parity across languages due to identical parallel FFT depth. \\
\textbf{Vocoder Fidelity} & HiFi-GAN produces crisp formants with minimal raspiness. & HiFi-GAN preserves vowel duration ratios accurately. & HiFi-GAN vocoding significantly outclasses Griffin-Lim. \\
\bottomrule
\end{tabularx}
\end{tcolorbox}
\vspace{-0.3cm}
\caption{Comparative evaluation matrix across linguistic and acoustic synthesis dimensions.}
\label{fig:comp_matrix}
\end{figure*}

\section{Reproducibility \& Open-Source Assets}

The entire codebase, model configurations, trained weights, generated audio samples, and variance plots have been open-sourced under the MIT License:

\begin{tcolorbox}[reportbox={Submission Metadata \& Repository}]
\small
\textbf{GitHub Repository}: \href{https://github.com/Tarak098/fastspeech2_project}{\texttt{https://github.com/Tarak098/fastspeech2\_project}} \\
\textbf{Environment}: Python 3.10+, PyTorch 2.1+, CUDA 12.1+ \\
\textbf{Dependencies}: \texttt{g2p\_en}, \texttt{soundfile}, \texttt{librosa}, \texttt{pyworld}, \texttt{pyyaml} \\
\textbf{Inference Execution}:
\begin{itemize}[leftmargin=*,noitemsep]
    \item English: \texttt{python inference.py config.yaml checkpoints/best\_model.pth}
    \item Telugu: \texttt{python inference.py config.yaml checkpoints/best\_model.pth}
\end{itemize}
\end{tcolorbox}

\section{Conclusion}
FastSpeech 2 demonstrates robust, non-autoregressive acoustic modeling across both stress-timed (English) and syllable-timed (Telugu) speech. Explicit conditioning on duration, pitch, and energy effectively eliminates the one-to-many ambiguity, providing stable phoneme alignment without attention failures or repetitive looping. The developed Telugu Unicode decomposition pipeline provides a generalizable framework for Indian Dravidian languages in modern neural TTS architectures.

\bibliographystyle{IEEEtran}
\begin{thebibliography}{9}
\bibitem{ren2020fastspeech2}
Y.~Ren, C.~Hu, X.~Tan, T.~Qin, S.~Zhao, Z.~Zhao, and T.-Y.~Liu, ``FastSpeech 2: Fast and high-quality end-to-end text to speech,'' in \emph{International Conference on Learning Representations (ICLR)}, 2021.

\bibitem{shen2018tacotron2}
J.~Shen \emph{et~al.}, ``Natural TTS synthesis by conditioning Wavenet on mel spectrogram predictions,'' in \emph{IEEE ICASSP}, 2018, pp. 4779--4783.

\bibitem{kong2020hifi}
J.~Kong, J.~Kim, and J.~Bae, ``HiFi-GAN: Generative adversarial networks for efficient and high fidelity speech synthesis,'' in \emph{NeurIPS}, 2020.

\bibitem{indictts}
A.~Baby \emph{et~al.}, ``Resources for Indian languages speech synthesis,'' in \emph{Proc. Interspeech}, 2016.
\end{thebibliography}

\end{document}
